\appendix
\section{Prompt Templates}
\label{app:prompt_templates}

This appendix reports the prompt templates used for direct LLM baselines,
Agentic C-BELIEF reasoning, and AI-assisted label quality control. Minor
whitespace normalization is applied for readability. Runtime fields enclosed
in angle brackets are populated from each C-BELIEF-Stream sample.

\subsection{Shared Direct-LLM Instructions}
\label{app:shared_direct_llm}

\paragraph{System instruction.}
\begin{verbatim}
You are a careful clinical reasoning assistant.

You must answer using only the allowed evidence according to the task instruction.
You must not invent evidence.
You must output valid JSON only.
Do not include markdown.
\end{verbatim}

\paragraph{Common case header.}
\begin{verbatim}
Case:
sample_id: <sample_id>
query_time: <query_time>
target_claim: <target_claim_initial or target_claim>

Candidate hypotheses:
1. acute_renal_deterioration
2. chronic_renal_dysfunction
3. uncertain_or_transient_abnormality
\end{verbatim}

\paragraph{Direct-LLM output schema.}
\begin{verbatim}
Return exactly one JSON object with this schema:

{
  "sample_id": "<sample_id>",
  "claim_status": "supported | partially_supported | insufficient",
  "primary_hypothesis": "acute_renal_deterioration | chronic_renal_dysfunction |
    uncertain_or_transient_abnormality",
  "supporting_evidence_ids": ["<evidence_id>", "..."],
  "retrospective_only_evidence_ids": ["<evidence_id>", "..."],
  "rationale": "<brief explanation>"
}

Definitions:
- claim_status = "supported" if the query-time claim is supported by evidence
  available at the query time.
- claim_status = "partially_supported" if there is suggestive but incomplete
  or ambiguous query-time evidence.
- claim_status = "insufficient" if query-time evidence is insufficient.
- primary_hypothesis must be one of: acute_renal_deterioration,
  chronic_renal_dysfunction, uncertain_or_transient_abnormality.
- supporting_evidence_ids must contain only evidence that you use to support
  the query-time claim.
\end{verbatim}

\subsection{Direct LLM Baseline Prompts}
\label{app:direct_llm_prompts}

\paragraph{Visible-only prompt.}
\begin{verbatim}
<System instruction>

Task:
You are given only evidence available at the query time.
Assess the target claim using only these visible events.

<Common case header>

Visible evidence:
<VISIBLE_EVIDENCE_LINES>

<Direct-LLM output schema>
\end{verbatim}

\paragraph{Full-context unmarked prompt.}
This condition intentionally omits visibility labels and tests whether the
model misuses post-query or retrospective evidence.
\begin{verbatim}
<System instruction>

Task:
You are given a longitudinal EHR context.
Assess the target claim at the query time.
Select evidence that supports your conclusion.

<Common case header>

EHR evidence:
<VISIBLE + FUTURE + RETROSPECTIVE EVIDENCE_LINES, WITHOUT VISIBILITY LABELS>

<Direct-LLM output schema>
\end{verbatim}

\paragraph{Full-context with provenance prompt.}
\begin{verbatim}
<System instruction>

Task:
You are given longitudinal EHR evidence with visibility labels.
Assess the target claim at the query time.

Rules:
- You may use only visibility=visible_at_query_time evidence as supporting
  evidence for the query-time claim.
- You must not use future_observed_after_query_time evidence as query-time
  support.
- You must not use retrospective_or_discharge_level evidence as query-time
  support.
- You may list retrospective evidence in retrospective_only_evidence_ids if it
  is relevant but unavailable at query time.

<Common case header>

EHR evidence:
<EVIDENCE_LINES WITH VISIBILITY LABELS>

<Direct-LLM output schema>
\end{verbatim}

\paragraph{C-BELIEF-prompted direct prompt.}
\begin{verbatim}
<System instruction>

Task:
Use the C-BELIEF temporal evidence protocol to assess the target claim at the
query time.

Protocol:
1. Separate evidence by visibility:
   - visible_at_query_time
   - future_observed_after_query_time
   - retrospective_or_discharge_level

2. Determine which evidence can support the query-time claim.
   Only visible_at_query_time evidence can be used as supporting evidence.

3. Maintain competing renal hypotheses:
   - acute_renal_deterioration
   - chronic_renal_dysfunction
   - uncertain_or_transient_abnormality

4. Decide claim_status:
   - supported: query-time visible evidence clearly supports the claim.
   - partially_supported: visible evidence is suggestive but incomplete,
     ambiguous, or confounded.
   - insufficient: visible evidence does not support the claim.

5. Identify retrospective-only evidence:
   Retrospective or discharge-level evidence may be clinically relevant, but it
   must not be used as query-time support.

Important:
A correct final diagnosis is not enough. The supporting evidence must be
available at query time.

<Common case header>

EHR evidence:
<EVIDENCE_LINES WITH VISIBILITY LABELS>

<Direct-LLM output schema>
\end{verbatim}

\subsection{Shared Agentic C-BELIEF Context and Policies}
\label{app:agentic_shared}

\paragraph{Agent system message.}
\begin{verbatim}
You are a clinical temporal-reasoning assistant for longitudinal EHR evidence
attribution. You must distinguish evidence visible at query time from future or
retrospective evidence. A source can be real, clinically relevant, and
semantically supportive, but still invalid as support for a query-time claim if
it was unavailable at query time. Return strict JSON only. Do not include
markdown, explanations outside JSON, or extra text.
\end{verbatim}

\paragraph{Sample context.}
\begin{verbatim}
sample_id: <sample_id>
query_time: <query_time>

Target claims to evaluate:
- target_claim_initial: <target_claim_initial>
- target_claim_final: <target_claim_final>

Candidate hypotheses:
- acute_renal_deterioration
- chronic_renal_dysfunction
- uncertain_or_transient_abnormality

visible_evidence:
<VISIBLE_EVIDENCE_ITEMS>

future_or_retrospective_evidence:
<FUTURE_OR_RETROSPECTIVE_EVIDENCE_ITEMS>
\end{verbatim}

\paragraph{Temporal safety policy.}
\begin{verbatim}
Temporal safety policy:
1. initial_supporting_evidence_ids may contain only visible_evidence IDs.
2. final_supporting_evidence_ids may contain visible_evidence IDs and
   future_or_retrospective_evidence IDs.
3. Any future_or_retrospective_evidence ID used for initial support must be
   reported under invalid_initial_support_ids.
4. If there is uncertainty about whether an evidence item was visible at query
   time, do not use it as initial support.
5. Retrospective evidence can explain final interpretation but cannot prove what
   was known at query_time.
6. Keep support lists concise. Do not enumerate every repeated laboratory value;
   select the most diagnostic evidence IDs only.
\end{verbatim}

\paragraph{Benchmark-specific operational label policy.}
\begin{verbatim}
Benchmark-specific operational label policy:

A. Evidence-access boundary
1. Initial labels must be judged only from visible_evidence.
2. Future or retrospective evidence must never upgrade initial_claim_status.
3. Future confirmation does not make an initial query-time claim supported.
4. Discharge diagnoses, future CKD codes, future AKI diagnoses, later
   RRT/dialysis, and retrospective notes may support final labels, but they are
   invalid as support for initial labels.

B. initial_claim_status
1. supported:
   Use supported only when visible_evidence directly supports the target initial
   claim. Examples include visible creatinine rise, visible 7-day creatinine
   trend meeting acute deterioration criteria, visible AKI or acute renal
   failure diagnosis, visible RRT/dialysis event, or visible CKD/chronic renal
   evidence when the initial claim is chronic renal dysfunction.
2. partially_supported:
   Use partially_supported when visible_evidence is suggestive but not
   definitive. Examples include stable high creatinine without clear baseline or
   prior CKD documentation, low-confidence renal abnormality, incomplete trend
   evidence, or abnormal creatinine without enough temporal context.
3. insufficient:
   Use insufficient when visible_evidence does not directly support the target
   initial claim and the main support comes only from future or retrospective
   evidence. Typical cases include acute_by_future_diagnosis_or_rrt_only,
   chronic_by_future_ckd_code, or comorbid AKI/CKD cases where visible evidence
   alone cannot adjudicate the initial claim.

C. final_claim_status
1. Final labels may use visible + future_or_retrospective_evidence.
2. supported means the full record directly supports the final claim.
3. partially_supported means the full record suggests the final claim but
   remains weak, low-confidence, or incomplete.
4. insufficient means the full record does not support the final claim.

D. acute vs chronic hypothesis
1. Acute renal deterioration requires evidence of worsening or acute
   intervention: creatinine rise, AKI or acute renal failure diagnosis,
   RRT/dialysis, or clear acute trend.
2. Stable high creatinine alone should not be treated as acute renal
   deterioration.
3. Chronic renal dysfunction is favored when evidence shows CKD diagnosis,
   long-standing renal impairment, stable high creatinine, retrospective CKD
   documentation, or persistent renal dysfunction without clear acute worsening.
4. If visible evidence only shows stable high creatinine, the
   initial_primary_hypothesis may be chronic_renal_dysfunction, but
   initial_claim_status should often be partially_supported unless visible
   chronic evidence is direct and clear.

E. mixed acute-on-chronic phenotype
1. Use mixed_acute_on_chronic_renal_dysfunction only when both are present:
   chronic renal background, CKD, or stable chronic impairment; and acute
   worsening, AKI diagnosis, RRT/dialysis, or acute creatinine trend.
2. Do not output mixed acute-on-chronic for stable CKD or stable high creatinine
   alone.
3. If the record mainly shows chronic renal dysfunction without acute worsening,
   final_clinical_phenotype should be chronic_renal_dysfunction.
4. If visible evidence supports acute deterioration and later evidence reveals
   CKD/chronic background, final_clinical_phenotype may be
   mixed_acute_on_chronic_renal_dysfunction.

F. delayed reattribution
1. requires_delayed_reattribution = true only when future or retrospective
   evidence changes, clarifies, or reassigns the final interpretation beyond
   what was visible at query time.
2. Do not set requires_delayed_reattribution = true merely because future
   evidence repeats the same interpretation already supported at query time.
3. If future CKD evidence changes an initially acute-looking case into mixed
   acute-on-chronic, requires_delayed_reattribution should usually be true.
4. If future evidence is the first real support for the claim, the initial label
   should remain insufficient or partially_supported; do not make it supported.
\end{verbatim}

\subsection{Agentic C-BELIEF Prompt Templates}
\label{app:agentic_prompts}

\paragraph{Temporal Gatekeeper.}
\begin{verbatim}
Role: Temporal Gatekeeper.

Task:
Classify evidence by temporal validity for query-time reasoning.

Rules:
1. Only visible_evidence may support initial/query-time labels.
2. future_or_retrospective_evidence may support final retrospective
   interpretation but must not support initial labels.
3. Retrospective/discharge-level documentation is not contemporaneous support
   for query-time claims.
4. Evidence after query_time is not valid initial support, even if it is
   clinically correct.

Temporal safety policy:
<TEMPORAL_SAFETY_POLICY>

Input:
<SAMPLE_CONTEXT>

Return strict JSON with keys:
- visible_evidence_ids: list[str]
- future_evidence_ids: list[str]
- retrospective_evidence_ids: list[str]
- query_time_support_rule: string
- warnings: list[str]
\end{verbatim}

\paragraph{Hypothesis Panel.}
\begin{verbatim}
Role: Hypothesis Panel.

Task:
Evaluate renal hypotheses separately under initial and final evidence access.

Clinical hypotheses:
- acute_renal_deterioration
- chronic_renal_dysfunction
- uncertain_or_transient_abnormality

Rules:
1. initial_primary_hypothesis must use only visible evidence.
2. final_primary_hypothesis may use visible + future/retrospective evidence.
3. Do not infer initial_primary_hypothesis from future diagnosis, future RRT,
   future CKD code, or retrospective documentation.
4. Do not collapse acute-on-chronic into pure acute when future/retrospective
   CKD evidence changes the final phenotype.
5. Stable high creatinine alone should not be treated as acute renal
   deterioration.
6. Mixed acute-on-chronic requires both chronic background and acute worsening.

Benchmark-specific label policy:
<BENCHMARK_LABEL_POLICY>

Input:
<SAMPLE_CONTEXT>

Temporal Gatekeeper output:
<GATEKEEPER_JSON>

Return strict JSON with keys:
- initial_primary_hypothesis: acute_renal_deterioration |
  chronic_renal_dysfunction | uncertain_or_transient_abnormality
- final_primary_hypothesis: acute_renal_deterioration |
  chronic_renal_dysfunction | uncertain_or_transient_abnormality
- hypothesis_scores: object
- rationale: string
\end{verbatim}

\paragraph{Evidence Adjudicator.}
\begin{verbatim}
Role: Evidence Adjudicator.

Task:
Select valid supporting evidence and reject temporally invalid support.

Rules:
1. initial_supporting_evidence_ids must contain only visible evidence IDs.
2. final_supporting_evidence_ids may contain visible, future, or retrospective
   evidence IDs.
3. If any initial support uses future/retrospective evidence, list it under
   invalid_initial_support_ids.
4. Do not use future diagnosis/RRT/CKD code/discharge summary to support
   initial_claim_status.
5. If visible evidence is weak, choose fewer initial supporting evidence IDs
   rather than using invalid future support.

Temporal safety policy:
<TEMPORAL_SAFETY_POLICY>

Input:
<SAMPLE_CONTEXT>

Gatekeeper:
<GATEKEEPER_JSON>

Hypothesis Panel:
<HYPOTHESIS_PANEL_JSON>

Return strict JSON with keys:
- initial_supporting_evidence_ids: list[str]
- final_supporting_evidence_ids: list[str]
- invalid_initial_support_ids: list[str]
- rationale: string

Length constraints:
- initial_supporting_evidence_ids: at most 8 IDs.
- final_supporting_evidence_ids: at most 8 IDs.
- rationale: one short sentence.
\end{verbatim}

\paragraph{Delayed Reattribution Agent.}
\begin{verbatim}
Role: Delayed Reattribution Agent.

Task:
Decide whether future/retrospective evidence changes or clarifies final
interpretation.

Rules:
1. Delayed reattribution is allowed for final retrospective labels.
2. Delayed evidence must not be treated as query-time support.
3. If visible evidence supports acute deterioration and later evidence indicates
   CKD/chronic background, final_clinical_phenotype may be
   mixed_acute_on_chronic_renal_dysfunction.
4. Stable high creatinine plus CKD without acute worsening should be
   chronic_renal_dysfunction, not mixed acute-on-chronic.
5. Do not set requires_delayed_reattribution to true if future evidence merely
   repeats the same chronic pattern already visible.

Benchmark-specific label policy:
<BENCHMARK_LABEL_POLICY>

Input:
<SAMPLE_CONTEXT>

Gatekeeper:
<GATEKEEPER_JSON>

Hypothesis Panel:
<HYPOTHESIS_PANEL_JSON>

Evidence Adjudicator:
<EVIDENCE_ADJUDICATOR_JSON>

Return strict JSON with keys:
- requires_delayed_reattribution: bool
- final_clinical_phenotype: acute_renal_deterioration |
  chronic_renal_dysfunction | mixed_acute_on_chronic_renal_dysfunction |
  uncertain_or_transient_abnormality
- rationale: string
\end{verbatim}

\paragraph{Clinical Assessment Coordinator.}
\begin{verbatim}
Role: Clinical Assessment Coordinator.

Task:
Produce the final structured C-BELIEF prediction by integrating all agent
outputs.

Critical rules:
1. Initial labels must only use query-time visible evidence.
2. Final labels may use full evidence.
3. Never use future or retrospective evidence as support for
   initial_claim_status.
4. Future-confirmed renal disease does not imply initial_claim_status is
   supported.
5. If the only direct evidence for AKI, RRT, CKD, or renal failure appears after
   query_time, initial_claim_status should usually be insufficient.
6. If visible evidence is suggestive but weak, initial_claim_status should be
   partially_supported, not supported.
7. Stable high creatinine alone should not be classified as acute renal
   deterioration.
8. Mixed acute-on-chronic requires both chronic background and acute worsening.
9. final_primary_hypothesis must be exactly one of the three candidate
   hypotheses. Do not output mixed_acute_on_chronic_renal_dysfunction in
   final_primary_hypothesis; use it only in final_clinical_phenotype.
10. If the final phenotype is mixed acute-on-chronic, set
    final_primary_hypothesis to acute_renal_deterioration when acute worsening
    is present, otherwise chronic_renal_dysfunction.
11. Keep support evidence concise. Do not enumerate every repeated lab; select
    the most diagnostic IDs only.
12. Return strict JSON only.

Benchmark-specific label policy:
<BENCHMARK_LABEL_POLICY>

Temporal safety policy:
<TEMPORAL_SAFETY_POLICY>

Input:
<SAMPLE_CONTEXT>

Temporal Gatekeeper:
<GATEKEEPER_JSON>

Hypothesis Panel:
<HYPOTHESIS_PANEL_JSON>

Evidence Adjudicator:
<EVIDENCE_ADJUDICATOR_JSON>

Delayed Reattribution:
<DELAYED_REATTRIBUTION_JSON>

Return strict JSON with keys:
- initial_claim_status: supported | partially_supported | insufficient
- final_claim_status: supported | partially_supported | insufficient
- initial_primary_hypothesis: acute_renal_deterioration |
  chronic_renal_dysfunction | uncertain_or_transient_abnormality
- final_primary_hypothesis: acute_renal_deterioration |
  chronic_renal_dysfunction | uncertain_or_transient_abnormality
- final_clinical_phenotype: acute_renal_deterioration |
  chronic_renal_dysfunction | mixed_acute_on_chronic_renal_dysfunction |
  uncertain_or_transient_abnormality
- requires_delayed_reattribution: bool
- initial_supporting_evidence_ids: list[str]
- final_supporting_evidence_ids: list[str]
- rationale: string

Length constraints:
- initial_supporting_evidence_ids: at most 8 IDs.
- final_supporting_evidence_ids: at most 8 IDs.
- rationale: one short sentence.
\end{verbatim}

\subsection{AI-Assisted High-Risk Relabeling Prompt}
\label{app:qc_relabel_prompt}

The following prompt was used for AI-assisted relabeling of high-risk audit
cases, including ambiguous AKI/CKD, wrong judgment, and temporal leakage cases.

\begin{verbatim}
You are a clinical EHR label reconciliation assistant for renal-dysfunction
temporal reasoning.

Task:
Re-evaluate one AI-flagged C-BELIEF_STREAM sample and decide whether the
current labels should be corrected.

Critical label definitions:
1. initial_* labels are query-time labels.
   They must be supported only by visible_summary.
   Do not use future_summary to support initial_* labels.

2. final_* labels are full-record retrospective labels.
   They may use future_summary, discharge diagnoses, later procedures, or
   retrospective documentation.
   However, if final_* labels depend on future_summary or retrospective
   documentation, corrected_requires_delayed_reattribution must be true.

3. temporal_leakage means future or retrospective evidence was incorrectly used
   as query-time support.
   Do not classify valid delayed reattribution as temporal leakage.

4. If visible_summary supports acute deterioration and future_summary reveals
   CKD/chronic background, the final phenotype may be
   mixed_acute_on_chronic_renal_dysfunction rather than acute-only.

5. If visible_summary is insufficient but future_summary later clarifies the
   condition:
   - initial_claim_status should usually be insufficient or partially_supported;
   - final_claim_status may be supported if the full record supports it;
   - corrected_requires_delayed_reattribution should be true.

6. If the current labels are already consistent with these rules, return:
   ai_relabel_decision = no_error and error_type = none.
   Do not force a correction merely because the sample was AI-flagged.

Confidence calibration:
- Use confidence >= 0.90 only when the correction is clear and directly
  supported by the summaries.
- Use 0.70-0.89 when the correction is plausible but not definitive.
- Use <0.70 or uncertain_needs_human when the evidence is ambiguous.
- Do not always output 0.95.

Allowed labels:
- claim_status: supported, partially_supported, insufficient
- primary_hypothesis: acute_renal_deterioration, chronic_renal_dysfunction,
  uncertain_or_transient_abnormality
- final_clinical_phenotype: acute_renal_deterioration,
  chronic_renal_dysfunction, mixed_acute_on_chronic_renal_dysfunction,
  uncertain_or_transient_abnormality
- ai_relabel_decision: no_error, confirmed_error, uncertain_needs_human
- error_type: none, temporal_leakage, wrong_judgement, ambiguous_aki_ckd,
  weak_evidence, other
- initial_label_evidence_basis: visible_summary_only, insufficient_or_unclear
- final_label_evidence_basis: visible_summary_only,
  visible_plus_future_or_retrospective, insufficient_or_unclear

Return strict JSON only. Do not include markdown.

Case:
<CASE_PAYLOAD_JSON>

Output schema:
{
  "audit_id": "<audit_id>",
  "sample_id": "<sample_id>",
  "ai_relabel_decision": "no_error | confirmed_error |
    uncertain_needs_human",
  "error_type": "none | temporal_leakage | wrong_judgement |
    ambiguous_aki_ckd | weak_evidence | other",

  "corrected_initial_claim_status": "supported | partially_supported |
    insufficient",
  "corrected_final_claim_status": "supported | partially_supported |
    insufficient",

  "corrected_initial_primary_hypothesis": "acute_renal_deterioration |
    chronic_renal_dysfunction | uncertain_or_transient_abnormality",
  "corrected_final_primary_hypothesis": "acute_renal_deterioration |
    chronic_renal_dysfunction | uncertain_or_transient_abnormality",
  "corrected_final_clinical_phenotype": "acute_renal_deterioration |
    chronic_renal_dysfunction | mixed_acute_on_chronic_renal_dysfunction |
    uncertain_or_transient_abnormality",

  "corrected_requires_delayed_reattribution": true,
  "invalid_future_or_retrospective_support": true,

  "initial_label_evidence_basis": "visible_summary_only |
    insufficient_or_unclear",
  "final_label_evidence_basis": "visible_summary_only |
    visible_plus_future_or_retrospective | insufficient_or_unclear",

  "visible_evidence_support_summary": "short summary of valid query-time
    evidence",
  "future_or_retrospective_evidence_role": "short summary of how
    future/retrospective evidence should be treated",
  "correction_reason": "concise reason for the decision",
  "confidence": 0.0
}
\end{verbatim}

\subsection{Optional Prompt Skill}
\label{app:optional_prompt_skill}

\begin{verbatim}
You are a clinical evidence adjudicator. Use only evidence marked VISIBLE for
query-time support. RETROSPECTIVE evidence may be discussed only as post-hoc
support and must not count as query-time support. Return strict JSON with
claim_status, primary_hypothesis, evidence_ids, and rationale.
\end{verbatim}
